% \documentclass[UTF8]{ctexart}
\documentclass{MMCStyle}
\usepackage[minted]{MyTemplate}
% \usepackage{siunitx}

% \title{\inline{LeeDo} 组2026年MathorCup论文撰写手册}
% \author{}
% \date{\today}

\begin{document}
% \maketitle
\begin{center}
\LARGE MathorCup论文撰写手册？\bigskip
\end{center}

\inline{DevlinNul} 自己写的规范，各位看官看看就好。

如果论文是多人协作的话，能用 \inline{git} 的团队就尽量用吧。

导言区需要包含的命令见附录。

参考文献的编译请用 bibtex , biber 也许可以通过。

\additionalSection{编码规范}

\begin{enumerate}
    \item 正文标点一律采用中文标点（全角），数学模式\textbf{中}的标点用英文标点。如\\
      $\forall a, b \in\R$ ，有 $a^2 + b^2 > 0$ 。
    \item 中文文本与标点之间不加空格，标点之间不加额外空格。
    \item 中文、英文、数字之间用一个空格隔开，如\\
      \inline{这是 1 段英文 English 示例。论文的英文是 article 。}\\
      数字是否用数学模式不作要求。
    \item 命令与文档文本之间用一个空格隔开，内联到正文中的无参数命令需补上空花括号。（如 \verb|\LaTeX| 应该写成 \verb|\LaTeX{}|）
    \item 与括号相连的文本无需空格隔开。如 \inline{括号（文本 Text 不 need）空格} 。
    \item 内联数学模式（行内公式）需与正文用空格隔开，如 \inline{正 $\int$ 文} 。
    \item 缩进不作要求。
    \item 表格原则上只允许使用全框表和三线表。
\end{enumerate}


\additionalSection{自定义的命令}

\begin{enumerate}
    \item 自然常数 $\e$ 需要使用直立体，可以使用命令 \verb|\e| 或者 \verb|\ee| 得到。
    \item 虚数单位 $\ii$ 需要使用直立体，可以使用命令 \verb|\ii| 得到。
    \item 物理量的单位使用 \inline{siuintx} 宏包得到，可用命令有 \verb|\num| 、 \verb|\unit| 和 \verb|\SI| ，不要使用 \verb|\qty| 。
    \item 可自动放缩的括号除了使用 \verb|\left(\right)| 外，还可以使用 \verb|\lr()| 。不要使用 \inline{physics} 宏包的 \verb|\qty()| 命令。
\end{enumerate}
\begin{table}[!htbp]
\centering
\caption{自定义命令表，省略前缀反斜杠，连接词的数学模式是可选的}
\begin{tabularx}{0.8\textwidth}{lYc}
\toprule
\textbf{符号} & \textbf{含义} & \textbf{渲染效果} \\
\midrule
\lstinline|lcm| & 最小公倍数 & $\lcm$ \\
\lstinline|arcsinh| & 反双曲正弦 & $\arcsinh$ \\
\lstinline|arccosh| & 反双曲余弦 & $\arccosh$ \\
\lstinline|arctanh| & 反双曲正切 & $\arctanh$ \\
\lstinline|arccoth| & 反双曲余切 & $\arccoth$ \\
\lstinline|arcsech| & 反双曲正割 & $\arcsech$ \\
\lstinline|arccsch| & 反双曲余割 & $\arccsch$ \\
\lstinline|ld| & 以 2 为底的对数 & $\ld$ \\
\lstinline|Cdot| & 比 \lstinline|cdot| 略大的点 & $\Cdot$ \\
\lstinline|N| & 自然数集 & $\N$ \\
\lstinline|Z| & 整数集 & $\Z$ \\
\lstinline|Q| & 有理数集 & $\Q$ \\
\lstinline|R| & 实数集 & $\R$ \\
\lstinline|C| & 复数集 & $\C$ \\
\lstinline|doesNotExist| & 不存在 & $\doesNotExist$ \\
\lstinline|LHS| & 左边 & $\LHS$ \\
\lstinline|RHS| & 右边 & $\RHS$ \\
\lstinline|almostEverywhere| & 几乎处处 & $\almostEverywhere$ \\
\lstinline|almostSurely| & 几乎总是 & $\almostSurely$ \\
\lstinline|suchThat| & 使得 & $\suchThat$ \\
\lstinline|eg| & 例如 & $\eg$ \\
\lstinline|ie| & 即 & $\ie$ \\
\lstinline|withoutLossOfGenerality| & 不失一般性 & $\withoutLossOfGenerality$ \\
\lstinline|etal| & 以及其他人 & $\etal$ \\
\lstinline|resp| & respectively & $\resp$ \\
\lstinline|viz| & 即 & $\viz$ \\
\lstinline|aka| & also known as & $\aka$ \\
\lstinline|diff| & 微元符号（括号要自己加） & $\diff{}$ \\
\lstinline|perm{2}{1}| & 排列数 & $\perm{2}{1}$ \\
\lstinline|comb{2}{1}| & 组合数 & $\comb{2}{1}$ \\
\lstinline|ii| & 虚数单位 & $\ii$ \\
\lstinline|ee| & 自然常数 & $\ee$ \\
\lstinline|sgn{x}| & 符号函数 & $\sgn{x}$ \\
\lstinline|na| & Not Available & $\na$ \\
% \lstinline|additionalSection| & section* & $\additionalSection$ \\
% \lstinline|addtionalSubsection| & subsection* & $\addtionalSubsection$ \\
% \lstinline|additionalSubsubsection| & subsubsection* &  \\
\bottomrule
\end{tabularx}
\end{table}

注： \verb|\gcd| 可以直接用。

\begin{table}
\centering
\caption{常用命令，由 physics 宏包提供（省略前缀反斜杠）}
\begin{tabularx}{0.8\textwidth}{lcY}
\toprule
\textbf{命令} & \textbf{效果} & \textbf{示例} \\
\midrule
\lstinline|eval{x}_s| & $\eval{x}_s$ & $\eval{\dfrac{f(x)}{g(y)}}_{\substack{
  x = 0 \\
  y = 1
}}^{\mathrm{up}}$ \\
\lstinline|derivative{y}{x}| & $\derivative{y}{x}$ & $\displaystyle\dv[2]{}{x}\lr(y z)$ \\
\bottomrule
\end{tabularx}
\end{table}

\appendix

\section{Miscellaneous}

\inline{MyTemplate} 这个自定义宏包的开关和支持的完整命令请自行看 \inline{.sty} 文件，因为这是自用的宏包，并不打算发布的，所以没有 manual 。

\section{导言区}

% 只是需要写某一 Part 的话采用的导言区：
\begin{solution}[只需要写某一 Part 的话采用的导言区：]
\begin{minted}{latex}
\documentclass{MMCStyle}
\usepackage{MyTemplate}
\end{minted}
\end{solution}

\begin{solution}[完整导言区：]
\begin{minted}{latex}
\documentclass{MMCStyle}
\usepackage{MyTemplate}

\bianhao{123456}
\tihao{ABCDEFGHI}
\timu{论文标题}
\keyword{content；content；content}
% \bibliographystyle{gbt7714-numerical}
\addbibresource{ref.bib}
\end{minted}
\end{solution}

\end{document}

